% Complete Chapter 7 introduction, original PDF 91-92 / printed 81-82.
% block: finite-element-forms-chapter-heading page:91 type:chapter
\clearpage\setcounter{section}{6}
\section{{有限元微分形式}}\label{chap:finite-element-differential-forms}

% block: finite-element-forms-hodge-problem page:91 type:ordinary-paragraph
{现在考虑 Lipschitz 多面体区域 $\Omega\subset\mathbb R^n$, 以及它的 $L^2$ de Rham 复形 \eqref{eq:l2-manifold-de-rham-complex}, 定义域复形为 \eqref{eq:manifold-de-rham-domain-complex}. 对任意 $k$, 片段}
\begin{equation}\label{eq:finite-element-forms-short-complex}
H\Lambda^{k-1}(\Omega)\xrightarrow{d}H\Lambda^k(\Omega)\xrightarrow{d}H\Lambda^{k+1}(\Omega)
\end{equation}
{确定了一个涉及两个外微分及其伴随的 Hodge Laplace 边值问题, 如 Section~\ref{sec:abstract-hodge-laplacian} 所述. 在 Theorem~\ref{thm:exterior-codifferential-adjoints} 中, 已将 $(d,H\Lambda^k)$ 的伴随确定为 $(\delta,\mathring H^*\Lambda^k)$, 因而 $k$-形式 Hodge Laplace 边值问题的强表述为}
\begin{equation}\label{eq:finite-element-forms-hodge-strong}
\begin{gathered}
(\delta d+d\delta)u=f-P_{\mathsf H}^kf\text{ in }\Omega,\\
\operatorname{tr}\star u=0,\ \operatorname{tr}\star du=0\text{ on }\partial\Omega,\qquad u\perp\mathsf H^k,
\end{gathered}
\end{equation}
{其中 $\mathsf H^k$ 为 $\Omega$ 上的调和 $k$-形式, 由微分方程 $dp=0$, $\delta p=0$ 和边界条件 $\operatorname{tr}\star p=0$ 确定. 在三维情形中, 已在 Section~\ref{sec:hodge-r3} 解释了此问题及相关问题. 在 Chapter~\ref{chap:approximation-hilbert-complexes} 中, 导出了使用混合表述和 Galerkin 方法成功离散化此问题所需的条件. 为此, 需要构造 $H\Lambda^{k-1}(\Omega)$ 的 Galerkin 子空间 $V_h^{k-1}$, 以及 $H\Lambda^k(\Omega)$ 的子空间 $V_h^k$. 除能够高效实现外, 这些子空间的关键要求是提供良好的逼近, 构成子复形, 并具有有界上链投影. 在这种情况下, 如 Chapter~\ref{chap:approximation-hilbert-complexes} 所示, 混合 Galerkin 方法相容且稳定, 并且得到最优阶收敛.}

% block: finite-element-forms-construction-goal page:91 type:ordinary-paragraph
{本章将通过构造有限元空间, 得到可高效实现的 $H\Lambda^k$ 的 Galerkin 子空间. 随后使用足够细的网格和/或足够高次的分片多项式形函数, 就能得到所需的逼近程度. 挑战在于保证这些空间构成具有有界上链投影的子复形.}

% block: complete-trimmed-families pages:91,92 type:ordinary-paragraph
{具体地, 给定 $\Omega$ 的任意单纯剖分 $\mathcal T_h$ (维数 $n\geq1$ 任意), 任意形式次数 $0\leq k\leq n$, 以及任意多项式次数 $r\geq1$, 将构造 $H\Lambda^k(\Omega)$ 的两个有限元子空间, 记为}
\[
\mathcal P_r\Lambda^k(\mathcal T_h)\quad{\text{和}}\quad\mathcal P_r^-\Lambda^k(\mathcal T_h).
\]
{前一个空间以次数不超过 $r$ 的多项式 $k$-形式的\emph{完全}空间作为形函数空间, 后一个则使用稍小的多项式空间, 按照 \BookCite{29}, 称其为\emph{修剪}空间. 随后给出几种选取两个此类空间的方法, 一个用于 $(k-1)$-形式, 另一个用于 $k$-形式, 以求解 Hodge Laplacian. 这些完全和修剪微分形式有限元空间族, 可看作离散化 de Rham 复形最典范的空间族. 这是 FEEC 的重要结论之一.}

% block: finite-element-forms-chapter-outline page:92 type:ordinary-paragraph
{本章安排如下. 在 Section~\ref{sec:finite-element-differential-forms} 定义完全空间和修剪空间 $\mathcal P_r\Lambda^k(\mathcal T_h)$ 与 $\mathcal P_r^-\Lambda^k(\mathcal T_h)$ 之前, 先在 Section~\ref{sec:complete-polynomial-forms}--\ref{sec:trimmed-polynomial-forms} 建立多项式系数微分形式的一些性质. 随后在 Section~\ref{sec:finite-element-space-properties}, 说明如何将这些空间组合为满足 Chapter~\ref{chap:approximation-hilbert-complexes} 误差分析所需假设的复形, 从而为我们的方法建立完整的误差界. 在 Section~\ref{sec:cubical-elements}, 简要说明基于长方体网格而非单纯形网格的有限元的类似结果. 最后在 Section~\ref{sec:historical-antecedents}, 将本章描述的单元与文献中的各种有限元联系起来.}

% block: finite-element-forms-notation page:92 type:ordinary-paragraph
{继续之前, 简要说明记号. 下一节将定义的有限元空间 $\mathcal P_r\Lambda^k(\mathcal T_h)$ 和 $\mathcal P_r^-\Lambda^k(\mathcal T_h)$, 由定义在区域 $\Omega$ 上的分片多项式微分 $k$-形式构成. 将对这些空间给出统一构造, 但每个空间都依赖于多项式次数 $r\geq0$, 介于 $0$ 与 $n=\dim\Omega$ 之间的形式次数 $k$, 以及单纯网格 $\mathcal T_h$. 还将用记号 $\mathcal P_r\Lambda^k$ 和 $\mathcal P_r^-\Lambda^k$ 表示次数为 $r$ 的完全和修剪多项式 (而非分片多项式) 微分 $k$-形式空间, 它们用于在每个单纯形上构造有限元空间. 这些多项式空间对每个 $n$, 每个 $r\geq0$ 及 $0\leq k\leq n$ 都有定义. 最后, 用 $\mathcal P_r\Lambda^k(T)$ 和 $\mathcal P_r^-\Lambda^k(T)$ 表示这些多项式限制在单个单纯形 $T$ 上的空间. 因此, 此记号类似于次数为 $r$ 的标准 Lagrange 有限元常用记号: $\mathcal P_r(\mathcal T_h)$ 表示有限元空间, $\mathcal P_r$ 表示次数为 $r$ 的多项式空间, $\mathcal P_r(T)$ 表示它们在一个单元上的限制.}
